%2multibyte Version: 5.50.0.2960 CodePage: 65001
%\newtheorem{definition}[definicion]{Definici�n}
%\newtheorem{example}[ejemplo]{Ejemplo}


\documentclass[9pt]{beamer}
%%%%%%%%%%%%%%%%%%%%%%%%%%%%%%%%%%%%%%%%%%%%%%%%%%%%%%%%%%%%%%%%%%%%%%%%%%%%%%%%%%%%%%%%%%%%%%%%%%%%%%%%%%%%%%%%%%%%%%%%%%%%%%%%%%%%%%%%%%%%%%%%%%%%%%%%%%%%%%%%%%%%%%%%%%%%%%%%%%%%%%%%%%%%%%%%%%%%%%%%%%%%%%%%%%%%%%%%%%%%%%%%%%%%%%%%%%%%%%%%%%%%%%%%%%%%
\usepackage{amsmath}
\usepackage{mathpazo}
\usepackage{hyperref}
\usepackage{multimedia}
\usepackage{times}
\usepackage[T1]{fontenc}

\setcounter{MaxMatrixCols}{10}
%TCIDATA{OutputFilter=LATEX.DLL}
%TCIDATA{Version=5.50.0.2960}
%TCIDATA{Codepage=65001}
%TCIDATA{<META NAME="SaveForMode" CONTENT="1">}
%TCIDATA{BibliographyScheme=Manual}
%TCIDATA{LastRevised=Tuesday, May 16, 2023 20:45:46}
%TCIDATA{<META NAME="GraphicsSave" CONTENT="32">}

\newtheorem{definicion}{Definici\'{o}n}
\newtheorem{ejemplo}{Ejemplo}
\newtheorem{ejemplos}{Ejemplos}
\newenvironment{stepenumerate}{\begin{enumerate}[<+->]}{\end{enumerate}}
\newenvironment{stepitemize}{\begin{itemize}[<+->]}{\end{itemize} }
\newenvironment{stepenumeratewithalert}{\begin{enumerate}[<+-| alert@+>]}{\end{enumerate}}
\newenvironment{stepitemizewithalert}{\begin{itemize}[<+-| alert@+>]}{\end{itemize} }
\usetheme{Rochester}
\usecolortheme[RGB={64,63,152}]{structure}
\usefonttheme{professionalfonts}
\input{tcilatex}
\begin{document}

\title[MAT-016]{Principios Matem\'{a}ticos}
\subtitle{Unidad III. Ecuaciones de Primer y Segundo Grado y sus Aplicaciones\\
Clase: Ecuaciones de primer grado}
\author[lgonzalez@santotomas.cl]{Profesor Luis Tulio Gonz\'{a}lez Alcaino \\
%EndAName
Magister en Matem\'{a}tica}
\institute[UST]{Universidad Santo Tom\'{a}s\\
Departamento Ciencias B\'{a}sicas - Talca }
\date{}
\maketitle

%TCIMACRO{\TeXButton{BeginFrame}{\begin{frame}}}%
%BeginExpansion
\begin{frame}%
%EndExpansion

\QTR{frametitle}{Contenidos de la clase}

\begin{itemize}
\item Ecuaciones de primer grado

\item Ecuaciones con coeficiente entero

\item Ecuaciones con coeficientes fraccionarios

\item Ejercicios
\end{itemize}

%TCIMACRO{\TeXButton{EndFrame}{\end{frame}}}%
%BeginExpansion
\end{frame}%
%EndExpansion

%TCIMACRO{\TeXButton{BeginFrame}{\begin{frame}}}%
%BeginExpansion
\begin{frame}%
%EndExpansion

\QTR{frametitle}{Ecuaciones}

\begin{definicion}
Una ecuaci\'{o}n es una igualdad entre dos expresiones separadas por el
signo igual, la cual presenta una o m\'{a}s variables, y que es verdadera
solamente para algunos valores de estas variables.
\end{definicion}

\pause

%TCIMACRO{\TeXButton{BeginBlok}{\begin{block}{Ejemplos de ecuaciones}}}%
%BeginExpansion
\begin{block}{Ejemplos de ecuaciones}%
%EndExpansion

\begin{stepenumeratewithalert}
\item $2x+5=3$ , ecuaci\'{o}n de primer grado, cuya soluci\'{o}n es $x=-1$,
en efecto, si reemplazamos $x=-1$ en la ecuaci\'{o}n, se tiene

$2\left( -1\right) +5=3$

$-2+5=3$

$3=3$

Con lo cual se obtiene la igualdad.

\item $\dfrac{2x}{5}+\dfrac{x}{10}=1$, ecuaci\'{o}n de primer grado con
coeficientes fraccionarios, cuya soluci\'{o}n es $x=2$
\end{stepenumeratewithalert}

%TCIMACRO{\TeXButton{EndBlok}{\end{block}}}%
%BeginExpansion
\end{block}%
%EndExpansion

%TCIMACRO{\TeXButton{EndFrame}{\end{frame}}}%
%BeginExpansion
\end{frame}%
%EndExpansion

%TCIMACRO{\TeXButton{BeginFrame}{\begin{frame}}}%
%BeginExpansion
\begin{frame}%
%EndExpansion

\QTR{frametitle}{Ecuaciones}

%TCIMACRO{\TeXButton{BeginBlok}{\begin{block}{Resolver ecuaciones}}}%
%BeginExpansion
\begin{block}{Resolver ecuaciones}%
%EndExpansion

\textbf{Resolver una ecuaci\'{o}n} significa encontrar todos los valores de
sus variables para los cuales la ecuaci\'{o}n es verdadera, (se cumple la
igualdad). Estos valores son llamados \textbf{soluciones de la ecuaci\'{o}n}
y se dice que satisfacen la ecuaci\'{o}n. Cuando s\'{o}lo est\'{a} implicada
una variable, \textbf{una soluci\'{o}n tambi\'{e}n es llamada ra\'{\i}z}. A
una letra o variable que representa una cantidad desconocida en una ecuaci%
\'{o}n se le denomina inc\'{o}gnita.

%TCIMACRO{\TeXButton{EndBlok}{\end{block}}}%
%BeginExpansion
\end{block}%
%EndExpansion

%TCIMACRO{\TeXButton{EndFrame}{\end{frame}}}%
%BeginExpansion
\end{frame}%
%EndExpansion

%TCIMACRO{\TeXButton{BeginFrame}{\begin{frame}}}%
%BeginExpansion
\begin{frame}%
%EndExpansion

\QTR{frametitle}{Ecuaciones}

%TCIMACRO{%
%\TeXButton{BeginBlok}{\begin{block}{Propiedades de igualdad de ecuaciones}}}%
%BeginExpansion
\begin{block}{Propiedades de igualdad de ecuaciones}%
%EndExpansion

Para resolver una ecuaci\'{o}n debemos tene presente las siguientes
propiedades de la igualdad:

\begin{stepitemizewithalert}
\item Al sumar o restar la misma cantidad en ambos lados de una ecuaci\'{o}%
n, la igualdad se mantiene.

\pause

\textbf{Ejemplo.} $x-3=5$ \ \ / $+3$

$\qquad \qquad x-3+3=5+3$

$\qquad \qquad x=8$

\item Al multiplicar o dividir por una cantidad distinta de cero en ambos
lados de una ecuaci\'{o}n, la igualdad se mantiene.

\pause

\textbf{Ejemplo.} $3x=12$/ $\div 3$

$\qquad \qquad \dfrac{3x}{3}=\dfrac{12}{3}$

$\qquad \qquad x=4$
\end{stepitemizewithalert}

%TCIMACRO{\TeXButton{EndBlok}{\end{block}}}%
%BeginExpansion
\end{block}%
%EndExpansion

%TCIMACRO{\TeXButton{EndFrame}{\end{frame}}}%
%BeginExpansion
\end{frame}%
%EndExpansion

%TCIMACRO{\TeXButton{BeginFrame}{\begin{frame}}}%
%BeginExpansion
\begin{frame}%
%EndExpansion

\QTR{frametitle}{Ecuaciones de primer grado en una variable (ecuaci\'{o}n
lineal)}

Para resolver una ecuaci\'{o}n de primer grado o lineal en una variable
realizamos operaciones en ella hasta obtener una ecuaci\'{o}n equivalente
cuya soluci\'{o}n es obvia. Esto significa una ecuaci\'{o}n en la que la
variable o inc\'{o}gnita queda aislada en un lado de la ecuaci\'{o}n.\pause

\begin{ejemplo}
Resolver la ecuaci\'{o}n $5x-6=3x+4$

\medskip \pause

\textbf{Soluci\'{o}n:} Empezamos por dejar los t\'{e}rminos que implican a $%
x $ en un lado y las constantes en el otro.

$5x-6=3x+4$\pause

$5x-3x=6+4$ (sumando $-3x$ y 6 a ambos miembros),\pause

$2x=10$ (simplificando),\pause

$x=\frac{10}{2}$ (dividiendo ambos miembros entre 2),\pause

$x=5\Longrightarrow $ La soluci\'{o}n de la ecuaci\'{o}n es 5.
\end{ejemplo}

%TCIMACRO{\TeXButton{EndFrame}{\end{frame}}}%
%BeginExpansion
\end{frame}%
%EndExpansion

%TCIMACRO{\TeXButton{BeginFrame}{\begin{frame}}}%
%BeginExpansion
\begin{frame}%
%EndExpansion

\bigskip \QTR{frametitle}{Ejemplos de ecuaciones de primer grado con
coeficientes enteros }

Resolver las siguientes ecuaciones:

\begin{stepenumerate}
\item $3-2x=5x+17$

\begin{stepitemize}
\item $x=-2$
\end{stepitemize}

\item $2\left[ (3x+1)-2(x+4)\right] -(3x+5)=5+x$

\begin{stepitemize}
\item $x=-12$
\end{stepitemize}

\item $2(x-1)=3(2-3x)-5(x-2)$

\begin{stepitemize}
\item $x=\frac{9}{8}$
\end{stepitemize}
\end{stepenumerate}

%TCIMACRO{\TeXButton{EndFrame}{\end{frame}}}%
%BeginExpansion
\end{frame}%
%EndExpansion

%TCIMACRO{\TeXButton{BeginFrame}{\begin{frame}}}%
%BeginExpansion
\begin{frame}%
%EndExpansion

\QTR{frametitle}{Ecuaciones de primer grado fraccionarias}

Para resolver ecuaciones fraccionarias de primer grado, se sugiere utilizar
la t\'{e}cnica de multiplicar por el m\'{\i}nimo com\'{u}n m\'{u}ltiplo de
los denominadores para dejar todos los coeficientes enteros. Luego resolver
una ecuaci\'{o}n resultante con coeficientes enteros.\pause

\begin{ejemplo}
Resolver $\dfrac{7x}{2}-\dfrac{9x}{4}$ $=\dfrac{2x-1}{4}+\dfrac{1}{2}$%
\pause

\textbf{Soluci\'{o}n:} Primero eliminamos las fracciones, multiplicando
ambos lados de la ecuaci\'{o}n por m\'{\i}nimo com\'{u}n m\'{u}ltiplo de los
denominadores, que es 4.

\begin{stepitemize}
\item $4\cdot \left( \dfrac{7x}{2}\right) -4\cdot \dfrac{9x}{4}=4\cdot
\left( \dfrac{2x-1}{4}\right) +4\cdot \left( \dfrac{1}{2}\right) $

\item $14x-9x=2x-1+2$

\item $5x=2x+1$

\item $5x-2x=1$

\item $3x=1$

\item $x=\dfrac{1}{3}\Longrightarrow $ \ La soluci\'{o}n o ra\'{\i}z de la
ecuaci\'{o}n es $x=\dfrac{1}{3}$
\end{stepitemize}
\end{ejemplo}

%TCIMACRO{\TeXButton{EndFrame}{\end{frame}}}%
%BeginExpansion
\end{frame}%
%EndExpansion

%TCIMACRO{\TeXButton{BeginFrame}{\begin{frame}}}%
%BeginExpansion
\begin{frame}%
%EndExpansion

\QTR{frametitle}{Ecuaciones de primer grado fraccionarias}

\begin{ejemplo}
Resolver $\dfrac{x-1}{3}+\dfrac{x}{2}$ $=\dfrac{x+3}{6}-2$

\textbf{Soluci\'{o}n: }Eliminamos las fracciones, multiplicando ambos lados
de la ecuaci\'{o}n por m\'{\i}nimo com\'{u}n m\'{u}ltiplo de los
denominadores, que es $6$.

\begin{stepitemize}
\item $6\cdot \left( \dfrac{x-1}{3}\right) +6\cdot \left( \dfrac{x}{2}%
\right) $ $=6\cdot \left( \dfrac{x+3}{6}\right) -6\cdot 2$

\item $2(x-1)+3x$ $=x+3-12$

\item $2x-2+3x=x-9$

\item $5x-2=x-9$

\item $5x-x=-9+2$

\item $4x=-7$

\item $x=-\dfrac{7}{4}$ $\ \Longrightarrow $ $\ $La soluci\'{o}n de la ecuaci%
\'{o}n es $x=-\dfrac{7}{4}$
\end{stepitemize}
\end{ejemplo}

%TCIMACRO{\TeXButton{EndFrame}{\end{frame}}}%
%BeginExpansion
\end{frame}%
%EndExpansion

%TCIMACRO{\TeXButton{BeginFrame}{\begin{frame}}}%
%BeginExpansion
\begin{frame}%
%EndExpansion

\QTR{frametitle}{Ejemplos de ecuaciones de primer fraccionarias}

Resolver las siguientes ecuaciones:

\begin{stepenumerate}
\item $\dfrac{x}{2}+\dfrac{2x}{5}=\dfrac{6x}{4}-\dfrac{5x-1}{10}$

\begin{stepitemize}
\item $x=-1$
\end{stepitemize}

\item $\dfrac{x-1}{2}-\dfrac{x+2}{4}=\dfrac{x+3}{6}$

\begin{stepitemize}
\item $x=18$
\end{stepitemize}
\end{stepenumerate}

%TCIMACRO{\TeXButton{EndFrame}{\end{frame}}}%
%BeginExpansion
\end{frame}%
%EndExpansion

%TCIMACRO{\TeXButton{BeginFrame}{\begin{frame}}}%
%BeginExpansion
\begin{frame}%
%EndExpansion

\QTR{frametitle}{Ejercicios}

\bigskip

Resolver para $x$ las siguientes ecuaciones:

\begin{enumerate}
\item $7x-3(2-3x)=4(2x+3)$

\item $\dfrac{x-3}{4}-\dfrac{x}{3}=\dfrac{x-1}{3}-\dfrac{x}{2}$

\item $\dfrac{5}{4}x+\dfrac{2}{5}x+\dfrac{3}{5}=\dfrac{1}{2}x+\dfrac{13}{10}%
x $

\item $4-(x+3)(x-1)=7-(x-2)^{2}$

\item $\dfrac{3-x}{4}+\dfrac{2x-1}{5}+1=\dfrac{x+1}{3}+\dfrac{x+3}{2}+\dfrac{%
2}{5}$

\item Si $4(3x+3)=5(6+2x)$, entonces el valor de $x$ es:

\item \textquestiondown Cu\'{a}l debe ser el valor de x para que la expresi%
\'{o}n \textquotedblleft\ $-2x$ \textquotedblright\ sea igual a la expresi%
\'{o}n \textquotedblleft\ $x-(8+x)$ \textquotedblright ?

\item \textquestiondown\ Cu\'{a}l es el valor de la inc\'{o}gnita $x$ para
que se cumpla que: \textquotedblleft\ $\frac{2}{3}x$ menos 3 sea igual a $2x$
m\'{a}s 1 \textquotedblright ?
\end{enumerate}

%TCIMACRO{\TeXButton{EndFrame}{\end{frame}}}%
%BeginExpansion
\end{frame}%
%EndExpansion

%TCIMACRO{\TeXButton{BeginFrame}{\begin{frame}}}%
%BeginExpansion
\begin{frame}%
%EndExpansion

\QTR{frametitle}{Soluciones de los ejercicios}

\begin{enumerate}
\item $\frac{9}{4}$

\item $5$

\item $4$

\item $\frac{2}{3}$

\item $-1$

\item $9$

\item $4$

\item $-3$
\end{enumerate}

%TCIMACRO{\TeXButton{EndFrame}{\end{frame}}}%
%BeginExpansion
\end{frame}%
%EndExpansion

\end{document}
